\chapter{Điều Người Lớn Hay Làm Học Trò Đau}
\markboth{Điều Người Lớn Hay Làm Học Trò Đau}{What Adults Often Do That Hurts Students}

\section{So Sánh Cho Nhanh}
\begin{truyen}{So Sánh Cho Nhanh}{Comparing Children Too Easily}
\chuhoa{C}{ó những câu người lớn nói rất quen: “Sao con không được như bạn ấy?”}
\textit{(There are familiar adult lines like: “Why can’t you be like that student?”)}

Nói ra thì nhanh, nhưng đứa trẻ nghe vào thường chỉ thấy mình bé đi.
\textit{(They are quick to say, but children often hear them as a shrinking of the self.)}

So sánh có thể kéo một đứa trẻ đi lên trong chốc lát, nhưng rất hay để lại gốc tự ti lâu hơn.
\textit{(Comparison may push a child briefly, but often leaves a deeper root of insecurity.)}

Không đứa trẻ nào lớn đẹp bằng cách phải liên tục mượn khuôn mặt của người khác.
\textit{(No child grows beautifully by constantly borrowing someone else’s face.)}

\end{truyen}

\begin{baihoc}
So sánh dễ tạo chuyển động, nhưng hiếm khi tạo được bình an để lớn bền.
\textit{(Comparison can create movement, but rarely creates the peace needed for lasting growth.)}
\end{baihoc}

\begin{ghinhoanh}
\textbf{Short English to remember:}

Comparison may create movement,

but it rarely creates peace for lasting growth.

\end{ghinhoanh}

\begin{tuhocnhanh}
\textbf{Gợi ý để dạy và tự nghĩ / Teaching and reflection prompts}

1. Vì sao so sánh làm trẻ đau? \textit{Why does comparison hurt children?}

2. Điều gì bị tạo ra sau câu nói này? \textit{What gets created after such a sentence?}

3. Có cách nào nhắc mà không so sánh? \textit{Is there a way to correct without comparing?}
\end{tuhocnhanh}
\ngancach

\section{Đùa Một Câu Nhưng Ở Lại Rất Lâu}
\begin{truyen}{Đùa Một Câu Nhưng Ở Lại Rất Lâu}{One Joke Can Stay for a Long Time}
\chuhoa{C}{ó câu người lớn nói đùa cho vui, nhưng học sinh mang theo nhiều tháng.}
\textit{(There are jokes adults say casually, yet students carry them for months.)}

Một câu chê ngoại hình, một câu gọi ngu, một câu gán tính cách tưởng nhẹ mà rất nặng.
\textit{(A joke about appearance, intelligence, or personality may sound light but land heavily.)}

Người lớn quên nhanh vì họ không phải người bị nói.
\textit{(Adults forget quickly because they are not the one who was spoken over.)}

Trẻ con có thể cười lúc đó và đau rất lâu sau đó.
\textit{(Children may laugh in the moment and hurt for a long time afterward.)}

\end{truyen}

\begin{baihoc}
Quyền lực trong lời nói của người lớn lớn hơn chính người lớn thường nghĩ.
\textit{(The power inside adult speech is greater than adults usually realize.)}
\end{baihoc}

\begin{ghinhoanh}
\textbf{Short English to remember:}

Adult words carry more power than adults often realize.

\end{ghinhoanh}

\begin{tuhocnhanh}
\textbf{Gợi ý để dạy và tự nghĩ / Teaching and reflection prompts}

1. Vì sao một câu đùa lại ở lâu? \textit{Why can one joke stay so long?}

2. Người lớn thường quên điều gì? \textit{What do adults often forget?}

3. Mình đang dùng lời nói như thế nào với học sinh? \textit{How are we using words with students?}
\end{tuhocnhanh}
\ngancach

\section{Chỉ Hỏi Điểm, Không Hỏi Người}
\begin{truyen}{Chỉ Hỏi Điểm, Không Hỏi Người}{Asking Only About Scores, Not the Person}
\chuhoa{N}{hiều đứa trẻ lớn lên trong cảm giác rằng chỉ khi làm tốt chúng mới đáng được hỏi han tử tế.}
\textit{(Many children grow up feeling that they are only worth warm attention when they perform well.)}

Người lớn hỏi điểm trước, mắng sau, rồi thôi.
\textit{(Adults ask about scores first, scold, and then stop there.)}

Lâu dần, học sinh cũng học cách nhìn chính mình bằng con số.
\textit{(Over time, students learn to view themselves through numbers as well.)}

Khi một đứa trẻ không được hỏi như một con người, chuyện học rất dễ biến thành chuyện chứng minh giá trị.
\textit{(When a child is not approached as a person, learning easily turns into a desperate attempt to prove worth.)}

\end{truyen}

\begin{baihoc}
Một câu hỏi về con người đôi khi chữa được nhiều hơn mười câu hỏi về kết quả.
\textit{(One question about the person can heal more than ten questions about performance.)}
\end{baihoc}

\begin{ghinhoanh}
\textbf{Short English to remember:}

One question about the person can heal more than ten about performance.

\end{ghinhoanh}

\begin{tuhocnhanh}
\textbf{Gợi ý để dạy và tự nghĩ / Teaching and reflection prompts}

1. Học sinh đang được hỏi như một kết quả hay như một con người? \textit{Is the student being approached as a result or as a person?}

2. Điều gì xảy ra khi trẻ tự nhìn mình bằng điểm? \textit{What happens when children view themselves through scores?}

3. Một câu hỏi khác có thể là gì? \textit{What could a different question sound like?}
\end{tuhocnhanh}
\ngancach

\section{Phán Nhanh Một Đứa Trẻ}
\begin{truyen}{Phán Nhanh Một Đứa Trẻ}{Judging a Child Too Quickly}
\chuhoa{C}{ó khi người lớn chỉ nhìn vài biểu hiện rồi kết luận rất gọn: lười, hỗn, vô ý thức, bất trị.}
\textit{(Sometimes adults see a few behaviors and quickly conclude: lazy, rude, careless, hopeless.)}

Kết luận nhanh giúp người lớn đỡ mất công nghĩ thêm.
\textit{(Fast judgment saves adults the effort of thinking further.)}

Nhưng một đứa trẻ thường phức tạp hơn rất nhiều so với cái nhãn đầu tiên.
\textit{(But a child is usually far more complex than the first label.)}

Phán nhanh là cách người lớn đơn giản hóa một con người còn đang lớn.
\textit{(Quick judgment is a way adults oversimplify a person who is still growing.)}

\end{truyen}

\begin{baihoc}
Muốn dạy đúng, người lớn phải chậm hơn một chút trong việc kết luận.
\textit{(If adults want to teach well, they need to be slower in drawing conclusions.)}
\end{baihoc}

\begin{ghinhoanh}
\textbf{Short English to remember:}

To teach well, adults must be slower in drawing conclusions.

\end{ghinhoanh}

\begin{tuhocnhanh}
\textbf{Gợi ý để dạy và tự nghĩ / Teaching and reflection prompts}

1. Vì sao người lớn hay phán nhanh? \textit{Why do adults judge too quickly?}

2. Cái giá của việc đó là gì? \textit{What is the cost of that habit?}

3. Mình có thể chậm lại ở bước nào? \textit{Where can we slow down?}
\end{tuhocnhanh}
\ngancach

\section{So Sánh Để Khích Lệ Nhưng Thành Ra Làm Đau}
\begin{truyen}{So Sánh Để Khích Lệ Nhưng Thành Ra Làm Đau}{Comparing to Motivate but Ending Up Hurting}
\chuhoa{N}{gười lớn rất hay nghĩ so sánh sẽ làm trẻ tỉnh ra. “Con nhìn bạn A đi.” “Em xem bạn B kìa.”}
\textit{(Adults often think comparison will wake children up. “Look at student A.” “See student B.”)}

Nhưng nhiều đứa trẻ không nghe được động lực trong câu đó.
\textit{(But many children do not hear motivation in such sentences.)}

Các em chỉ nghe thấy rằng bản thân mình đang kém đi trong mắt người lớn.
\textit{(They only hear that they are becoming smaller in adult eyes.)}

So sánh có thể đẩy trẻ chạy. Nhưng cũng có thể làm trẻ xấu hổ và rút vào trong.
\textit{(Comparison can make a child run. It can also make a child feel ashamed and retreat inward.)}

\end{truyen}

\begin{baihoc}
Không phải cách thúc nào cũng nuôi được ý chí. Có cách chỉ nuôi mặc cảm.
\textit{(Not every push builds willpower. Some only feed insecurity.)}
\end{baihoc}

\begin{ghinhoanh}
\textbf{Short English to remember:}

Not every push builds willpower.

Some only feed insecurity.

\end{ghinhoanh}

\begin{tuhocnhanh}
\textbf{Gợi ý để dạy và tự nghĩ / Teaching and reflection prompts}

1. Vì sao so sánh dễ làm đau? \textit{Why does comparison hurt so easily?}

2. Học sinh nghe câu so sánh theo cách nào? \textit{How do students usually hear comparison?}

3. Mình có thể thay bằng cách khích lệ nào khác? \textit{What can replace comparison-based motivation?}
\end{tuhocnhanh}
\ngancach

\section{Mắng Ở Chỗ Đông Người}
\begin{truyen}{Mắng Ở Chỗ Đông Người}{Scolding in Front of a Crowd}
\chuhoa{C}{ó người lớn nghĩ mắng trước lớp cho em chừa. Nhưng có những đứa trẻ không nhớ bài học nào cả sau chuyện đó. Các em chỉ nhớ mình đã xấu hổ thế nào.}
\textit{(Some adults think public scolding will teach a lesson. But some children do not remember the lesson at all. They only remember the humiliation.)}

Lòng tự trọng của học sinh ở tuổi này rất dễ trầy.
\textit{(A student’s self-respect at this age scratches easily.)}

Những vết xước vì bị bêu trước đám đông có thể âm ỉ lâu hơn một lỗi sai cụ thể.
\textit{(The wounds of public shame can last longer than the original mistake.)}

Dạy cho nhớ không nhất thiết phải làm cho nhục.
\textit{(Teaching a lasting lesson does not require making a child feel degraded.)}

\end{truyen}

\begin{baihoc}
Có những lỗi nên được sửa riêng, vì phẩm giá của trẻ cũng là một phần của giáo dục.
\textit{(Some mistakes should be corrected in private, because a child’s dignity is also part of education.)}
\end{baihoc}

\begin{ghinhoanh}
\textbf{Short English to remember:}

A child’s dignity is also part of education.

\end{ghinhoanh}

\begin{tuhocnhanh}
\textbf{Gợi ý để dạy và tự nghĩ / Teaching and reflection prompts}

1. Vì sao bêu trước đám đông dễ để lại vết sâu? \textit{Why does public shame leave a deep mark?}

2. Điều gì cần được giữ khi sửa lỗi? \textit{What must be protected while correcting mistakes?}

3. Mình có đang lẫn giữa nghiêm và làm nhục không? \textit{Are we confusing firmness with humiliation?}
\end{tuhocnhanh}
\ngancach

\section{Nghe Nửa Câu Rồi Kết Luận Cả Con Người}
\begin{truyen}{Nghe Nửa Câu Rồi Kết Luận Cả Con Người}{Hearing Half a Story and Judging the Whole Person}
\chuhoa{N}{gười lớn đôi khi rất nhanh: nghe một lời kể, thấy một hành vi, rồi gắn luôn một nhãn lên học sinh.}
\textit{(Adults can be very fast: they hear one report, see one behavior, and quickly place a label on a student.)}

Nhưng trẻ con cũng như người lớn, không phải lúc nào cũng tệ theo đúng hình người khác kể.
\textit{(But children, like adults, are not always as bad as the story makes them seem.)}

Một kết luận vội dễ biến thành cách đối xử vội.
\textit{(A rushed conclusion easily becomes rushed treatment.)}

Và một khi đã bị nhìn méo, học sinh rất khó có cơ hội được nhìn lại cho đúng.
\textit{(Once a student is seen in a distorted way, it becomes hard to be seen fairly again.)}

\end{truyen}

\begin{baihoc}
Trước khi sửa một học sinh, người lớn nhiều khi cần sửa sự vội của chính mình.
\textit{(Before correcting a student, adults often need to correct their own haste first.)}
\end{baihoc}

\begin{ghinhoanh}
\textbf{Short English to remember:}

Before correcting a student, adults often need to correct their own haste first.

\end{ghinhoanh}

\begin{tuhocnhanh}
\textbf{Gợi ý để dạy và tự nghĩ / Teaching and reflection prompts}

1. Vì sao người lớn hay kết luận nhanh? \textit{Why do adults rush to conclusions?}

2. Nhãn sai gây hại thế nào? \textit{How does a wrong label harm a student?}

3. Mình cần kiểm tra lại điều gì trước khi phán xét? \textit{What should we re-check before judging?}
\end{tuhocnhanh}
\ngancach

\section{Bắt Trẻ Xin Lỗi Khi Chưa Hiểu Mình Sai Ở Đâu}
\begin{truyen}{Bắt Trẻ Xin Lỗi Khi Chưa Hiểu Mình Sai Ở Đâu}{Forcing an Apology Before the Child Understands}
\chuhoa{C}{ó những lời xin lỗi được nói rất nhanh vì người lớn yêu cầu. Nhưng bên trong đứa trẻ chưa hiểu chuyện gì vừa xảy ra.}
\textit{(Some apologies are spoken very quickly because adults demand them, while inside the child does not yet understand what really happened.)}

Khi đó, lời xin lỗi chỉ còn là thủ tục để hết chuyện.
\textit{(In that case, the apology becomes only a procedure to end the incident.)}

Giáo dục không chỉ là dạy trẻ nói “em xin lỗi.”
\textit{(Education is not only teaching children to say “I’m sorry.”)}

Giáo dục còn là giúp trẻ thấy vì sao hành động của mình đã chạm vào người khác.
\textit{(It is also helping them see how their action touched another person.)}

\end{truyen}

\begin{baihoc}
Một lời xin lỗi có ý nghĩa phải đi sau một sự hiểu, không chỉ đi sau một mệnh lệnh.
\textit{(A meaningful apology should follow understanding, not only instruction.)}
\end{baihoc}

\begin{ghinhoanh}
\textbf{Short English to remember:}

A meaningful apology should follow understanding, not just instruction.

\end{ghinhoanh}

\begin{tuhocnhanh}
\textbf{Gợi ý để dạy và tự nghĩ / Teaching and reflection prompts}

1. Vì sao xin lỗi ép buộc dễ rỗng? \textit{Why do forced apologies often feel empty?}

2. Trẻ cần hiểu điều gì trước khi xin lỗi? \textit{What should a child understand before apologizing?}

3. Mình đang dạy lời hay đang dạy sự hiểu? \textit{Are we teaching words or understanding?}
\end{tuhocnhanh}
